\documentclass[10pt,a4paper]{amsart}
\usepackage[margin=19mm]{geometry}
\usepackage{amsmath,amssymb,mathtools}
\usepackage[scr=rsfs,cal=euler]{mathalfa}
\usepackage{xfrac}
\usepackage[unicode,hidelinks,bookmarks=false]{hyperref}
\newcommand{\ie}{i.e.\ }
\newcommand{\cf}{cf.\ }
\newcommand{\resp}{respectively\ }
\newcommand{\Subsection}{\subsection}
\providecommand{\nobreakdash}{\nobreak\hbox{-}}
\setlength{\emergencystretch}{2em}
\multlinegap0pt
\usepackage{amssymb}
\usepackage{tikz}
\newtheorem{corollary}{Corollary}
\newtheorem{lemma}{Lemma}
\newtheorem{definition}{Definition}
\newtheorem{observation}{Observation}
\title{Partitioning perfect graphs into comparability graphs}
\author{Andr\'as Gy\'arf\'as, M\'arton Marits, G\'eza T\'oth}
\date{Source: arXiv:2408.13523v2 (CC BY 4.0)}
\begin{document}
\maketitle

\noindent\emph{Source and licence.} Partitioning perfect graphs into comparability graphs. Version arXiv:2408.13523v2 (CC BY 4.0), distributed by arXiv under the Creative Commons Attribution 4.0 International licence, http://creativecommons.org/licenses/by/4.0/, as recorded in the arXiv licence metadata for this exact version and shown on the abstract page https://arxiv.org/abs/2408.13523v2. Full licence notice: \texttt{licenses/arxiv-2408.13523-CC-BY-4.0.txt}.\par\smallskip

\noindent\textit{Editorial scope.} Counted unit: the article's lemma rec (source0.tex lines 484--486). The article's complete original proof of it is delivered with it (source0.tex lines 488--598, one \texttt{proof} environment of 111 lines). No mathematics is written, repaired or re-derived: the statement and the proof are the article's own lines, in the article's own notation, and no retained line is substituted or re-flowed.  Retained uncounted as the source-local context the proof needs: lines 416--419.  Nothing was omitted from any retained span, so the package carries no omission record.  No retained line contains a citation, so the wrapper carries no bibliography and no bibliography entry.  No retained line contains a figure environment, an image-inclusion command or drawing code, so no image is delivered and the package declares no asset dependency.  Editorial formatting only, in addition: an AMS \texttt{amsart} preamble that restates the article's own declarations and macros used by the retained lines, namely the environments \texttt{corollary}, \texttt{lemma}, \texttt{definition}, \texttt{observation} on separate counters, together with the standard packages those lines need.  The retained environments print in this document as Corollary 1, Lemma 1, Definition 1, Observation 1 in order of appearance, whereas the article prints the same items with the numbering of its own full document.  The complete native source of the article as distributed in the arXiv e-print for this exact version is bundled unchanged as \texttt{sources/164213/source0.tex}.
\begin{corollary}\label{shift}
  $p(S_n)\le \lceil \log(\chi(S_n))\rceil = \lceil\log(\lceil\log(n)\rceil)\rceil$
  and the partition can be done by bipartite graphs.
\end{corollary}

\begin{lemma}\label{rec}  We have $p'(4)=1$. For $n\ge 3$ we have
 $$p'(n^2)\le  p'(n)+ \lceil\log(\lceil\log(n)\rceil)\rceil + 2.$$
\end{lemma}

\begin{proof}[Proof of Lemma \ref{rec}]
%\noindent{\bf Proof of Lemma \ref{rec}.}
It is easy to check that $p'(4)=1$, i.e. $F_4$ is a comparability graph.
Indeed, in the vertex ordering
$$(2,4),(1,2),(3,4),(2,3),(1,3),(1,4),$$

orienting all the ten edges of $F_4$ forward, we get a transitive orientation, thus $F_4$ is indeed a comparability graph.
Let $n\ge 3$ now.
Assume that
$P=H_1,\dots,H_{p'(n)}$ is a partition of the edges of $F_n$ into $p'(n)$
comparability subgraphs. Partition the base point set $[n^2]$ of $F_{n^2}$ into $n$ consecutive blocks of size
$n$,  $B_1,\dots,B_n$, where for $i\in [n]$,

$$B_i=\{(i-1)n+1,\dots,in\}.$$




The intervals (that is, vertices of
$F_{n^2}$)
with both endpoints in the same block $B_i$
are called {\em inner} intervals.
The intervals with endpoints in different blocks $B_i,B_j,
i\ne j$  are the {\em outer} intervals.

Let $e$ be the edge of $F_{n^2}$ defined by a pair of
outer intervals $[a, b]$ and $[c, d]$,
where $a\in B_{a'}$,  $b\in B_{b'}$,   $c\in B_{c'}$,   $d\in B_{d'}$.
Then $e$ is called {\em outer-nested, outer-crossing, outer-touching, or outer-degenerate}
if the intervals $[a',b'],[c',d']$ are nested, crossing, touching or coincide (i.e. $a'=c',b'=d'$).


For $i\in [n]$, let $F(B_i)$ denote the subgraph of $F_{n^2}$ induced by the inner intervals of $B_i$.

\begin{definition}
  The {\em blow-up of type (a)} of a graph $G$ is the following.
  Replace vertices  $v_i$ of $G$ by pairwise disjoint vertex sets $V_i$.
  For each edge $e=(v_i,v_j)$ of $G$, add a complete bipartite graph between
  the corresponding sets $V_i$ and $V_j$.

  The {\em blow-up of type (b)} of a graph $G$ is the following.
Replace vertices  $v_i$ of $G$ by pairwise disjoint vertex sets $V_i$.
  For each edge $e=(v_i,v_j)$ of $G$, add an arbitrary set of edges between
  the corresponding sets $V_i$ and $V_j$.
\end{definition}


The following observation is obvious.

\begin{observation}\label{blow}
  The blow-up of type (a) of a comparability graph is a comparability graph.
    The blow-up of type (b) of a bipartite graph is a bipartite graph.
\end{observation}

Returning to the proof of Lemma \ref{rec}, $F_{n^2}$ is partitioned as follows.


\begin{itemize}

\item (i)  Clearly, $F(B_i)$ is isomorphic to $F_n$ so it has a partition
  $P_i= H^i_1, \ldots, H^i_{p'(n)}$ into $p'(n)$ comparability graphs. Since unions of vertex
  disjoint comparability graphs are comparability graphs,
  $P^*= H^*_1, \ldots, H^*_{p'(n)}$
is a partition on the edges of $F_{n^2}$ defined by
inner intervals into comparability graphs $H^*_j=\cup_{i=1}^nH^i_j$.

Let $NC\subseteq F_{n^2}$ be a graph whose vertices correspond to the outer intervals and edges
  correspond to outer-nested and
  outer-crossing pairs of intervals.
  Clearly, the graph $NC$ is a blow-up of type (a) of $F_n$. Therefore, by Observation \ref{blow}, there is a partition
  of its edges into $p'(n)$ comparability graphs,
 $P'=H'_1,\dots,H'_{p'(n)}$.
Again, since  $P'$ is defined on
outer intervals while $P^*$ is on inner intervals and
   unions of vertex
   disjoint comparability graphs are comparability graphs, we can take the union of the partitions
   $P^*$ and $P'$ to get
   a partition of the {\em edges defined by inner intervals and by outer-nested and outer-crossing pairs} into
   $p'(n)$ comparability graphs.


 \item (ii) Consider now the graph $OS$ of {\em intersecting outer-touching intervals} of $F_{n^2}$.
   An outer-touching edge is a pair of outer-touching intervals $[i,j], [k,l]$ $i<j$,$k<l$ where
   $i\in B_a$, $j, k\in B_b$ and $a<b<c$, so it corresponds to a touching pair $[a, b], [b, c]$. Therefore,
   $OS$ is a blow-up of type (b) of $S_n$. Observe, that a pair of outer-touching intervals are not necessarily intersecting,
   that's why we have a blow-up of type (b) and not (a).



   Corollary \ref{shift} gives a partition of the edges of $S_n$ defined by touching pairs of intervals
   into $\lceil\log(\lceil\log(n)\rceil)\rceil$ bipartite graphs.
   By Observation \ref{blow}, it corresponds to a partition of the edges of $OS$
   into the same number of bipartite graphs. Since bipartite graphs are comparability graphs, we
   obtained a partition of $OS$ into  $\lceil\log(\lceil\log(n)\rceil)\rceil$ comparability graphs.

 \item (iii) Let $OD$ be the graph of {\em intersecting outer-degenerate intervals} of $F_{n^2}$.
   For any interval $[i, j]$, $1\le i<j\le n^2$, $i\in B_a$, $j\in B_b$, the type of $[i,j]$ is the pair $(a,b)$.
   Observe that two outer intervals form an intersecting outer-degenerate pair if and only if they are of the same type.
Therefore, the graph $OD$ is the disjoint union of ${n\choose 2}$ cliques, which is a comparability graph.





\item (iv) Finally, there are the edges of $F_{n^2}$ that correspond to a crossing or nested pair of an inner and an outer interval.
But these edges form a bipartite graph, so it is another comparability graph.

\end{itemize}

Since all edges of $F_{n^2}$ are partitioned in (i)-(iv), summing up, we get the statement of Lemma \ref{rec}.
\end{proof}

\end{document}
